%% kapa: exam=1 round=22 year=2011 period=2 form=A subject=accounting subject_ko=회계학 seq=40 number=40
%% source: Q-Net 감정평가사 기출문제 (cst00309) artlSeq=5036790 file=extracted/1008/5036790/2011년 제22회 감정평가사 1차시험 기출문제/2교시.hwp
%% conversion: kordoc 4.24.51 -> tools/gen1.py (auto); review=proofread-sonnet; flags=none
\begin{question}{40}
(주)서울은 단일 제품을 생산하고 있으며 20×1년 재공품에 관한 자료는 다음과 같다.

{\small\centering
\begin{tabular}{lr|lr}
\hline
기초재공품 수량 & - & 완성량 & 80개 \\
당기착수량 & \uline{130개} & 기말재공품 수량 & \uline{50개} \\
합계 & \uline{130개} & 합계 & \uline{130개} \\
\hline
\end{tabular}\par}

기말재공품의 가공원가 완성도는 40\%이다. (주)서울은 당기에 직접노무시간 660시간을 투입하였다. 회사의 제품 단위당 표준직접노무시간은 6시간이고, 표준임률은 ￦3,000이다. 당기에 실제로 발생된 직접노무원가가 ￦2,100,000이었다면, 직접노무원가의 능률차이는? \par
\begin{choices}
  \item ￦120,000 불리
  \item ￦180,000 불리
  \item ￦120,000 유리
  \item ￦180,000 유리
  \item ￦540,000 불리
\end{choices}
\end{question}
\answer{2}
